% ============================================================
\section{Detector Precision Utility is Context-Dependent}
\label{sec:mech}
% ============================================================

\noindent\textbf{Non-additive interactions.}
We measured all $\binom{12}{2}=66$ pairwise interactions on a 6-bit base
(all blocks at 6-bit; $U(\varnothing)$ has score\_drop $0.0298$).
Table~\ref{tab:interaction} summarizes: the average absolute centered
interaction is $0.00135$ versus mean $|\Delta(i\mid\varnothing)|=0.00171$,
a ratio of $0.79$.
Extremes are instructive: \texttt{BLK2.3}$+$\texttt{BLK2.5} are nearly
useless alone yet interact by $+0.0060$, while two of the best singletons
(\texttt{BLK2.0}/\texttt{BLK2.2}) underperform their sum by $-0.0047$.

\begin{table}[t]
\centering
\small
\caption{Pairwise interactions, 6-bit base (calib-64, GT-free).}
\label{tab:interaction}
\begin{tabular}{lr}
\toprule
statistic & value \\
\midrule
mean $I_{ij}$ & $+0.0005$ \\
std $I_{ij}$ & $0.0017$ \\
range & $[-0.0047, +0.0060]$ \\
mean $|I_{ij}|$ & $0.00135$ \\
mean $|\Delta(i\mid\varnothing)|$ & $0.00171$ \\
ratio & $0.79$ \\
\bottomrule
\end{tabular}
\end{table}

\begin{figure}[t]
\centering
\includegraphics[width=0.98\linewidth]{figures/fig1_reversal.pdf}
\caption{Singleton-vs-conditional ranking under growing context:
(a)~pairwise order reversal rate (all pairs vs.\ large-gap pairs,
$|{\rm gap}|>0.002$); (b)~top-1 mismatch rate.}
\label{fig:reversal}
\end{figure}

\noindent\textbf{Rank reversal.}
If interactions were small or homogeneous, singleton rankings would still
be usable. They are not: comparing, for every context $S$, the conditional
ranking $\Delta(i\mid S)$ against the singleton ranking
$\Delta(i\mid\varnothing)$ over the remaining candidates, $22.6\%$,
$24.7\%$, $26.4\%$ of pairwise orders reverse at $|S|{=}1,2,3$
(Fig.~\ref{fig:reversal}; $12/66/220$ contexts, top-1 mismatch
$25/61/62\%$). Reversals are mostly moderate in size, but the largest
reversals change decisions: in context
$S=\{\texttt{BLK2.5},\texttt{BLK3.0}\}$ conditional measurement prefers
\texttt{BLK2.3} over the singleton-favored \texttt{BLK2.4}, a $+0.0070$
utility swing ($\approx 4\times$ the mean singleton effect).

\smallskip
\noindent\textbf{Five-seed replication.}
Four additional 64-image pools give reversal rates of
$22.6$--$32.0\%$ (all $\geq 20\%$) and interaction-to-singleton ratios
of $0.50$--$1.12$ (mean $0.76$); per-pair values are seed-dependent, so
we report distribution-level statements, which all replicate.

\smallskip
\noindent\textbf{Decision impact.}
Utility-level differences translate to decisions. At an equal repair
budget (4 of 12 blocks at 8-bit, 4-bit base), selecting blocks by
singleton rankings of the same instrument gives $31.39$ AP---about
chance-level among repair sets ($31.07$)---while conditional measurement
gives $34.95$ ($+3.56$), and the oracle repair set gives $35.32$. Exact
enumeration of all $\binom{12}{4}=495$ sets confirms the ordering
(top-4: $41/495$ vs.\ $11/495$); predicted and
measured gains correlate at Spearman $0.72$--$1.00$ (legal signals)
vs.\ $-0.16$--$0.60$ (feature proxy).
